\documentclass{scrreprt}

\KOMAoptions{fontsize=11pt, paper=a4} % Schriftgröße und Papierformat setzen.
\KOMAoptions{DIV=16} % Parameter mit dem man den Seitenrand ändern kann.


\usepackage{mathtools}
\usepackage{amssymb}
\usepackage{hyperref}
\usepackage[all]{hypcap}
\usepackage{subcaption}
\usepackage{float}
\usepackage[german]{babel}
\usepackage{siunitx}
\usepackage{physics}
\usepackage[version=4]{mhchem}
\usepackage{pdfpages}
\usepackage{longtable}

\sisetup{
	separate-uncertainty = true,
	multi-part-units=brackets,
	number-unit-product = \text{ },
	mode=math,
	exponent-product = \cdot,
	output-product = \cdot,
	output-decimal-marker = {,}
}

\usepackage{ circuitikz }

\usepackage[backend=biber, style=numeric]{biblatex} % do not change this, the whole thing will break
% \usepackage{bibtex}
% \usepackage[backend=bibtex]{biblatex}

% \usepackage[headsepline=true,markcase=upper]{scrlayer-scrpage}
% \automark[chapter]{chapter}
% \ohead*{521: $\gamma$-Spektroskopie} % outer head on even & odd pages
% \setkomafont{pagehead}{\tiny\bfseries}




% common macros so i don't have to import packages or type too much
\def\diff#1#2{\frac{\text{d}#1}{\text{d}#2}}
\def\diffp#1#2{\frac{\partial#1}{\partial#2}}
\def\imply{\Longrightarrow}
\def\d#1{\text{d}#1\text{ }}
%\def\equiv{\Longleftrightarrow} % synonym for \iff


\def\todo#1{
	\textbf{\textcolor{blue}{[#1]}}
}
\def\moszwei{$\text{MoS}_2$}





\def\jafp#1#2{
	\begin{center}
		\includegraphics[width=0.8\textwidth]{#1}
		%\includegraphics[height=0.4\textheight]{#1}
		\captionof{figure}{#2}
		\label{fig:#1}
	\end{center}
}

\def\jafps#1#2#3{
	\begin{figure}[h]
		\centering
		\includegraphics[width=#3\textwidth]{#1}
		%\includegraphics[height=0.4\textheight]{#1}
		\captionof{figure}{#2}
		\label{fig:#1}
	\end{figure}
}

\def\jafpsw#1#2#3{
	\begin{figure*}[h]
		\centering
		\includegraphics[width=#3\textwidth]{#1}
		%\includegraphics[height=0.4\textheight]{#1}
		\captionof{figure}{#2}
		\label{fig:#1}
	\end{figure*}
}


\def\jafpp#1#2#3#4{
	\begin{figure}[H]
		\centering
		\begin{subfigure}{0.45\textwidth}
			\centering
			\includegraphics[width = \textwidth]{#1}
			\caption{}
			\label{fig:#1}
		\end{subfigure}
		\begin{subfigure}{0.45\textwidth}
			\centering
			\includegraphics[width = \textwidth]{#3}
			\caption{}
			\label{fig:#3}
		\end{subfigure}
		\caption{a) #2 und b) #4}
	\end{figure}
}

\def\jafpsh#1#2#3{
	\begin{figure}[H]
		\centering
		\includegraphics[width=#3\textwidth]{#1}
		%\includegraphics[height=0.4\textheight]{#1}
		\captionof{figure}{#2}
		\label{fig:#1}
	\end{figure}
}




\def\ergo{\Rightarrow} %silly stuff for me

\addtokomafont{chapterprefix}{\raggedleft}
\addtokomafont{chapter}{\Large}
\addtokomafont{section}{\large}
\addtokomafont{subsection}{\large}
\addtokomafont{subsubsection}{\large}


\makeatletter
\renewcommand\chapter{\thispagestyle{plain}%
  \global\@topnum\z@
  \@afterindentfalse
  \secdef\@chapter\@schapter}
\makeatother

\DeclareSIUnit\angstrom{\protect \text  {Å}}


\usepackage{graphicx}


%\titlehead{\centering\includegraphics[width=0.6\textwidth]{mos2_usb.jpg}} %for a title picture
\graphicspath{{../figs/}}
\author{Leonie Dessau \& Lena Beckmann}

\setlength\parindent{0pt}
